\section{Operation Object}

An Operation object represents a machining operation within a setup.

Operations define the machining action performed on the workpiece and reference the resources required to execute that action, such as tools, toolpaths, coordinate systems, offsets, compensation, notes, and process parameters.

Each operation is stored as an independent JSON object within the \texttt{operations} directory of a setup.

Example location:

\begin{verbatim}
setups/set-1/operations/op-2.json
\end{verbatim}

\subsection{Operation Structure}

An Operation object typically contains:

\begin{itemize}
\item an operation identifier
\item an optional descriptive name
\item a strategy identifier
\item artifact references to associated resources
\end{itemize}

The operation role and machining kind are defined by the setup operation index entry rather than repeated within the operation object.

Example:

\begin{verbatim}
{
  "id": "op-2",
  "name": "Face1",
  "strategy": "face",

  "artifacts": [
    {
      "role": "tool",
      "kind": "json",
      "path": "../../../tools/main/2.json",
      "present": true
    },
    {
      "role": "toolpath",
      "kind": "apt",
      "path": "cldata/op-2.apt",
      "present": true
    },
    {
      "role": "operation_csys",
      "kind": "json",
      "path": "csys/op-2.json",
      "present": true
    }
  ]
}
\end{verbatim}

\subsection{Operation Identifier}

Each operation must define a unique identifier using the \texttt{id} field.

Example:

\begin{verbatim}
"id": "op-2"
\end{verbatim}

The identifier allows the operation to be indexed by the setup and referenced by other objects within the package.

\subsection{Strategy}

The \texttt{strategy} field identifies the machining strategy used by the operation.

Examples include:

\begin{itemize}
\item \texttt{face}
\item \texttt{contour}
\item \texttt{adaptive}
\item \texttt{drill}
\end{itemize}

The strategy field is informational and may vary depending on the exporting CAM system.

\subsection{Operation Artifacts}

Operations reference associated resources through artifacts.

Common artifacts associated with operations include:

\begin{itemize}
\item tool definitions
\item toolpath files
\item coordinate system definitions
\item cutting condition definitions
\item compensation definitions
\item operation offsets
\item notes or documentation
\item extensions and addons
\end{itemize}

These resources are stored as separate files within the operation directory.

Software implementations may interpret only the artifacts they support while ignoring unsupported elements.
